% Integration-ready research section; labels bil:*; bibliography keys below.
\section{A Gamma normalization for independently shifted Hurwitz tails}
\label{bil:section}

The diagonal square law in the preceding report leaves open what happens
when both tail orders and both shifts vary independently. A normalization
by an elementary Gamma quotient gives a finite endpoint formula in all
four variables. The resulting weight is determined by the two orders;
an arbitrary additional outer power is not being evaluated here.

\subsection{The normalization and its difference equation}

For $\Re z>0$ write
\begin{align}
 \mathcal B_s(z)
 &=2^{-s}\left\{\zeta\left(s,\frac z2\right)
                    -\zeta\left(s,\frac{z+1}{2}\right)\right\},
                    \label{bil:alternating-hurwitz}\\
 \ell(z)&=\frac12\log2+\log\Gamma\left(\frac{z+1}{2}\right)
                     -\log\Gamma\left(\frac z2\right),\qquad
 h(z)=e^{\ell(z)}.\label{bil:h}
\end{align}
The difference in \eqref{bil:alternating-hurwitz} has its removable
value at $s=1$ and is entire in $s$. All logarithms of Gamma in
\eqref{bil:h} are the holomorphic branches on the right half-plane that
are real on the positive axis. Powers of $h$ mean $e^{s\ell}$.
The elementary recurrences give
\begin{equation}\label{bil:basic-recurrences}
 \mathcal B_s(z)+\mathcal B_s(z+1)=z^{-s},\qquad
 \ell(z)+\ell(z+1)=\log z,\qquad h(z)h(z+1)=z.
\end{equation}
The alternating Hurwitz recurrence and its order derivatives are
classical; see Hu--Kim \cite{bil:HuKim}. The quotient $h$ is also
classical: in the notation of Wang--Hu--Kim \cite{bil:WangHuKim},
$h(z)=\sqrt{\pi/2}/\widetilde\Gamma(z)$.

Set
\begin{equation}\label{bil:tau-v}
 \tau_s(z)=h(z)^s\mathcal B_s(z),\qquad
 V_s(z)=\frac{z^s-(z-1)^s}{h(z)^{2s}}\quad(\Re z>1).
\end{equation}
Neither $h$ nor this rewriting of $\mathcal B_s$ is asserted to be a
new special function. Their useful feature is the homogeneous equation
\begin{equation}\label{bil:homogeneous}
 \boxed{\tau_s(z-1)-\tau_s(z+1)=V_s(z)\tau_s(z).}
\end{equation}
Indeed, multiplication of the two adjacent recurrences for
$\mathcal B_s$ by the appropriate powers of $h$ gives
\[
 \tau_s(z\pm1)=h(z)^{-s}
                  -\left(\frac{h(z\pm1)}{h(z)}\right)^s\tau_s(z).
\]
Subtracting and using \eqref{bil:basic-recurrences} proves
\eqref{bil:homogeneous}.

\subsection{Two exact bilinear endpoint laws}

\begin{theorem}[Independent orders and shifts]\label{bil:master}
Let $\Re a,\Re b>0$. For $p,q\in\mathbb C$ put
\[
 u_n=\tau_p(n+a+1),\qquad v_n=\tau_q(n+b+1),\qquad
 C_n=u_nv_{n+1}+u_{n+1}v_n,\quad
 D_n=u_nv_{n+1}-u_{n+1}v_n.
\]
These definitions include $n=-1$. For every integer $N\ge0$,
\begin{align}
 &\sum_{n=0}^{N}
 \{V_p(n+a+1)+V_q(n+b+1)\}u_nv_n
       =C_{-1}-C_N,\label{bil:finite-plus}\\
 &\sum_{n=0}^{N}(-1)^n
 \{V_p(n+a+1)-V_q(n+b+1)\}u_nv_n
       =D_{-1}+(-1)^ND_N.\label{bil:finite-minus}
\end{align}
If $\Re(p+q)>0$, both series converge absolutely and locally normally,
and hence
\begin{align}
 &\sum_{n\ge0}\{V_p(n+a+1)+V_q(n+b+1)\}u_nv_n\notag\\
 &\hspace{8mm}=\tau_p(a)\tau_q(b+1)+\tau_p(a+1)\tau_q(b),
                                      \label{bil:plus}\\
 &\sum_{n\ge0}(-1)^n\{V_p(n+a+1)-V_q(n+b+1)\}u_nv_n\notag\\
 &\hspace{8mm}=\tau_p(a)\tau_q(b+1)-\tau_p(a+1)\tau_q(b).
                                      \label{bil:minus}
\end{align}
All fixed mixed derivatives in $p,q,a,b$ may be taken term by term
in this domain. The second equality also holds as an ordinary,
possibly conditional, series on $\Re(p+q)>-2$.
\end{theorem}

\begin{proof}
The homogeneous equation yields
\[
 C_{n-1}-C_n=\{V_p(n+a+1)+V_q(n+b+1)\}u_nv_n,
\]
and
\[
 D_{n-1}+D_n=\{V_p(n+a+1)-V_q(n+b+1)\}u_nv_n.
\]
The first identity telescopes. Multiplying the second by $(-1)^n$
also telescopes, since consecutive signs are opposite. This proves
both finite identities, including their lower signs.

The standard Hurwitz expansion and the Gamma quotient expansion give,
uniformly on compact sets of orders and shifts,
\begin{align}
 \mathcal B_s(z)&=\frac12z^{-s}
                       \left(1+\frac{s}{2z}+O(z^{-2})\right),\notag\\
 h(z)&=z^{1/2}\left(1-\frac1{4z}+O(z^{-2})\right),\notag\\
 \tau_s(z)&=\frac12z^{-s/2}
                       \left(1+\frac{s}{4z}+O(z^{-2})\right),
                         \label{bil:asymptotics}\\
 V_s(z)&=\frac{s}{z}+O(z^{-2}).\notag
\end{align}
These expansions hold for arbitrary bounded complex $s$, with
sufficiently many terms in the Hurwitz expansion when needed. Thus
the summands in \eqref{bil:plus} and \eqref{bil:minus} are
$O(n^{-1-\Re(p+q)/2})$, locally uniformly, and
$C_N=O(N^{-\Re(p+q)/2})$. The leading terms cancel in $D_N$, giving
$D_N=O(N^{-1-\Re(p+q)/2})$. This proves absolute normal convergence
and the endpoint limits when $\Re(p+q)>0$. Cauchy's formula gives
the derivative assertion. Finally the finite identity
\eqref{bil:finite-minus} itself proves ordinary convergence on
$\Re(p+q)>-2$, since its remaining endpoint tends to zero there.
\end{proof}

The Hurwitz and Gamma expansions invoked here are classical
\cite{bil:DLMF}; the contribution is the bilinear finite endpoint
calculus built from them. On setting $p=q=s$ and $a=b$, division
of \eqref{bil:plus} by two gives precisely
\[
 \sum_{n\ge0}\{(n+a+1)^s-(n+a)^s\}\mathcal B_s(n+a+1)^2
 =\mathcal B_s(a+1)-a^s\mathcal B_s(a+1)^2.
\]
Thus the earlier square law is a diagonal specialization of this
four-variable family.

\begin{corollary}[The first boundary contact]\label{bil:contact}
For $p+q=0$, the ordinary symmetric series in
\eqref{bil:finite-plus} converges absolutely and satisfies
\begin{align}
 &\sum_{n\ge0}\{V_p(n+a+1)+V_{-p}(n+b+1)\}
           \tau_p(n+a+1)\tau_{-p}(n+b+1)\notag\\
 &=\tau_p(a)\tau_{-p}(b+1)
        +\tau_p(a+1)\tau_{-p}(b)-\frac12.\label{bil:contact-half}
\end{align}
\end{corollary}
\begin{proof}
The $n^{-1}$ terms in the sum of the two $V$'s cancel, so the
summand is $O(n^{-2})$. Equation \eqref{bil:asymptotics} gives
$C_N\to1/2$. Now pass to the limit in \eqref{bil:finite-plus}.
\end{proof}
In particular the ordinary boundary value differs by $1/2$ from the
analytic endpoint continuation from $\Re(p+q)>0$. On the diagonal,
the division by two explains the earlier $1/4$ contact. The present
statement does not confuse those two notions of value.

\subsection{All independent Stieltjes orders}

Use the convention
\[
 \zeta(1+t,z)=\frac1t+
       \sum_{j\ge0}\frac{(-1)^j\gamma_j(z)}{j!}t^j.
\]
For $r\ge0$ define the following explicit finite expressions:
\begin{align}
 T_r(z)&=\frac{h(z)}2
  \sum_{j=0}^{r}\binom rj
  (\ell(z)-\log2)^{r-j}(-1)^j
  \left\{\gamma_j\left(\frac z2\right)
         -\gamma_j\left(\frac{z+1}{2}\right)\right\},
                                        \label{bil:T-jet}\\
 J_r(z)&=\frac1{h(z)^2}
  \left\{z(\log z-2\ell(z))^r
      -(z-1)(\log(z-1)-2\ell(z))^r\right\}.
                                        \label{bil:J-jet}
\end{align}
Here $\Re z>0$ suffices in \eqref{bil:T-jet}, and $\Re z>1$ is
used in \eqref{bil:J-jet}. Direct coefficient extraction gives
\[
 T_r(z)=\left.\partial_s^r\tau_s(z)\right|_{s=1},\qquad
 J_r(z)=\left.\partial_s^r V_s(z)\right|_{s=1}.
\]
For example $T_0(z)=h(z)\mathcal B_1(z)$ and $J_0(z)=h(z)^{-2}$.

\begin{theorem}[Every independent mixed Stieltjes jet]
\label{bil:stieltjes-all}
For $r,s\ge0$ and $\Re a,\Re b>0$, write $z_n=n+a+1$ and
$w_n=n+b+1$. Put
\begin{align*}
 X_{r,s}(z,w)&=\sum_{j=0}^{r}\binom rj
                         J_j(z)T_{r-j}(z)T_s(w),\\
 Y_{r,s}(z,w)&=\sum_{k=0}^{s}\binom sk
                         T_r(z)J_k(w)T_{s-k}(w).
\end{align*}
Then the following are ordinary absolutely convergent identities:
\begin{align}
 \sum_{n\ge0}\{X_{r,s}(z_n,w_n)+Y_{r,s}(z_n,w_n)\}
 &=T_r(a)T_s(b+1)+T_r(a+1)T_s(b),\label{bil:jets-plus}\\
 \sum_{n\ge0}(-1)^n\{X_{r,s}(z_n,w_n)-Y_{r,s}(z_n,w_n)\}
 &=T_r(a)T_s(b+1)-T_r(a+1)T_s(b).\label{bil:jets-minus}
\end{align}
\end{theorem}
\begin{proof}
Apply $\partial_p^r\partial_q^s$ to \eqref{bil:plus} and
\eqref{bil:minus} at $(p,q)=(1,1)$. The product rule gives the two
finite binomial sums exactly. Each individual scalar summand in these
sums is $O(n^{-2}\log^{r+s}(n+2))$, so no regularization or
cancellation between divergent scalar series is implicit.
\end{proof}

For the first mixed jet, $r=s=1$, the symmetric summand is explicitly
\[
 \{J_0(z)+J_0(w)\}T_1(z)T_1(w)
 +J_1(z)T_0(z)T_1(w)+J_1(w)T_1(z)T_0(w).
\]
Its sum is $T_1(a)T_1(b+1)+T_1(a+1)T_1(b)$. Thus independent
first Stieltjes coefficients at different shifts already occur in
an exact finite endpoint identity, and \eqref{bil:jets-plus}--
\eqref{bil:jets-minus} give every higher pair without an inference
from a single diagonal derivative.

\subsection{Anchored primitives and a pure Gamma square identity}

Differentiating \eqref{bil:h} shows
\begin{equation}\label{bil:h-derivative}
 \ell'(z)=\mathcal B_1(z),\qquad \tau_1(z)=h'(z).
\end{equation}
Hence the $(p,q)=(1,1)$ laws have elementary primitives in the Gamma
quotient itself. Let $A=[a_0,a_1]$ and $B=[b_0,b_1]$ be oriented
real intervals with positive endpoints, and define
\[
 U_A(t)=h(t+a_1)-h(t+a_0),\qquad
 R_A(t)=\frac1{h(t+a_0)}-\frac1{h(t+a_1)},
\]
with analogous notation for $B$. Double integration of
\eqref{bil:plus} and \eqref{bil:minus}, using
$\int h'/h^2=-1/h$, gives
\begin{align}
 \sum_{n\ge1}\{R_A(n)U_B(n)+U_A(n)R_B(n)\}
 &=U_A(0)U_B(1)+U_A(1)U_B(0),\label{bil:primitive-plus}\\
 \sum_{n\ge1}(-1)^{n-1}\{R_A(n)U_B(n)-U_A(n)R_B(n)\}
 &=U_A(0)U_B(1)-U_A(1)U_B(0).\label{bil:primitive-minus}
\end{align}
These anchored identities have no unspecified integration constants.
Uniform convergence on the compact rectangles justifies integration;
alternatively $U_A(n)=O(n^{-1/2})$ and $R_A(n)=O(n^{-3/2})$
give absolute convergence directly.

\begin{corollary}[A Gamma quotient increment square]\label{bil:gamma-square}
For positive real $a,b$,
\begin{equation}\label{bil:gamma-square-formula}
 \boxed{\sum_{n=1}^{\infty}
   \frac{\{h(n+b)-h(n+a)\}^{2}}{h(n+a)h(n+b)}
   =\{h(b)-h(a)\}\{h(b+1)-h(a+1)\}.}
\end{equation}
\end{corollary}
\begin{proof}
Take the same interval $A=B=[a,b]$ in
\eqref{bil:primitive-plus} and divide by two.
There is also a particularly short finite certificate. If
$U_n=h(n+b)-h(n+a)$, the recurrence $h(z)h(z+1)=z$ gives
\[
 U_{n-1}U_n-U_nU_{n+1}
       =\frac{U_n^2}{h(n+a)h(n+b)}.
\]
Thus the sum through $n=N$ equals $U_0U_1-U_NU_{N+1}$, and
$U_NU_{N+1}=O(N^{-1})$ tends to zero.
\end{proof}
At $(a,b)=(1,2)$ the right side is
\begin{equation}\label{bil:gamma-square-example}
 3-\frac\pi2-\frac4\pi.
\end{equation}
Indeed $h(1)=\sqrt{2/\pi}$, $h(2)=\sqrt{\pi/2}$, and
$h(3)=2\sqrt{2/\pi}$. Equation \eqref{bil:gamma-square-formula}
also extends holomorphically to $\Re a,\Re b>0$, with the algebraic
square as written and no complex conjugation.

\subsection{Scope and further targets}

The classical ingredients are the alternating Hurwitz recurrence,
Gamma shift formula, and a two-solution discrete Lagrange identity.
Products of zeta tails already have a broad literature, including
Hoffman's multiple-zeta treatment \cite{bil:Hoffman}. The claims here
are the displayed independently parameterized endpoint laws, their
explicit mixed Stieltjes coefficients, and the anchored Gamma
consequences, as additions to the inspected project material. A
literature search did not locate these particular forms; this is not
a proof of priority. No numerical independence assertion is used.

Two precise extensions remain. First, adjoining an independent outer
power to \eqref{bil:plus} generally breaks the displayed telescoper;
one should describe the additional coordinate and its moving polar
hyperplanes. Second, at $p+q=-2m$ with $m>0$, the exact upper endpoint
provides a systematic continuation scheme, but the complete
independent-order contact algebra is not established by
Corollary~\ref{bil:contact} alone.
